\begin{threeparttable}
\begin{tabular}{lccccccccc}
\toprule
 & \multicolumn{3}{c}{Exposure weighted} & \multicolumn{3}{c}{Unweighted} & & & \\
\cmidrule(lr){2-4} \cmidrule(lr){5-7}
Window & Mean & Std. dev. & IQR & Mean & Std. dev. & IQR & Nb. banks & ESS & W \\
\midrule
2009-2014 & 0.21    & 0.59 & 1.13 & $-0.21$ & 0.61 & 0.70 & 26 & 8.38  & 0.18 \\
2014-2019 & 0.14    & 1.53 & 3.05 & 0.79    & 1.05 & 1.83 & 42 & 3.39  & 0.51 \\
2019-2024 & $-0.35$ & 0.39 & 0.53 & $-0.62$ & 0.54 & 0.42 & 61 & 7.93  & 0.29 \\
\midrule
Pooled     & 0.00    & 1.00 & 1.17 & 0.00    & 1.00 & 1.04 & 79 & 16.53 & 0.17 \\
\bottomrule
\end{tabular}
\begin{tablenotes}[flushleft]
\footnotesize
\item Notes: The table describes the accumulated five-year bank supply shocks $\hat{\gamma}_{bt,5}$, estimated on the full sample rather than the leave-one-out counterparts used by the instrument, whose only difference is the exclusion of one market at a time. All statistics are computed on the estimation sample of the main specification ($1{,}390$ market-years; three five-year windows ending in 2014, 2019, and 2024).
\item Within each block (exposure-weighted and unweighted), shocks are standardized to zero mean and unit variance over the estimation panel using the corresponding weights, so the pooled means and standard deviations equal 0 and 1 by construction. IQR is the difference between the 75th and 25th percentiles.
\item The weight $w_{bt}$ is bank $b$'s share of baseline credit in each markets in window $t$, normalized to sum to one within each window for the window rows and across all bank-window cells for the pooled row. ESS $=1/\sum w_{bt}^{2}$ is the effective sample size of the shocks \citep{borusyak2022quasi}. W denotes the largest normalized weight.
\item Nb. banks counts banks with an estimated shock in each window; the
pooled row reports distinct banks (79, across 129 bank-window cells).
\end{tablenotes}
\end{threeparttable}